% =============================================================================
% 1.2 国内外研究现状
% 草稿版本：v0.1 (2026-05-05)
% 字数：约 2200 字
% 风格规则：v0.2 风格（无 \textbf / \textit / 引例括号 / 双引号 / 破折号）
% 引用：\cite{ref1, ref2, ref5, ref6, ref7, ref8, ref9, ref10, ref11, ref12,
%        ref13, ref15, ref16, ref17, ref18, ref19, ref20, ref21, ref22, ref23,
%        ref24, ref25, ref26, ref27, ref28, ref29, ref30}
% 依赖宏包：booktabs / tabularx（主稿已加载）
% =============================================================================

\subsection{国内外研究现状}
\label{sec:related-work}

本课题涉及的研究方向主要包括气象预报领域的 AI 应用、基于大语言
模型的 AI 智能体技术、检索增强生成与工具学习三个方向。本节按
这三个方向分别综述其研究现状与发展趋势，并在最后归纳现有研究
的不足以及本课题的定位。

\subsubsection*{1.2.1 \quad 气象预报领域的 AI 应用}

近年来深度学习在气象预报领域取得了突破性进展。华为公司的盘古
气象大模型 \cite{ref16} 首次实现了基于三维神经网络的中期全球
天气预报，在多项气象要素的预报精度上达到甚至超越欧洲中期天气
预报中心的传统数值预报模型。Google DeepMind 的 GraphCast
\cite{ref17} 通过图神经网络在 10 天全球天气预报任务上展现出
卓越性能。国内的风乌 \cite{ref26} 与伏羲 \cite{ref27} 模型分别
在中期预报精度与长时效级联预报方向取得重要进展。盛杰等 \cite{ref7}
对临近气象预报大模型风雷的 V1 版本进行了系统检验评估。李怡然等
\cite{ref6} 提出了基于气象大模型的门控时空融合网络闪电预报方法。
彭跃华等 \cite{ref8} 探索了盘古气象大模型与控制论的融合用于
台风预测分析。Rasp 等 \cite{ref25} 建立的 WeatherBench 2 为数据
驱动的气象模型提供了标准化评测基准。

上述工作主要聚焦于预报精度的提升，回答了如何预报得更准这一问题,
但较少关注预报数据的智能检索与分析，即如何让用户更便捷地获取并
理解预报信息。代刊等 \cite{ref5} 指出，大语言模型在天气预报中的
应用不应局限于预报本身，还应拓展到预报解释、数据检索、报告生成
等业务辅助环节。李特等 \cite{ref2} 进一步提出了 AI 智能体在天气
预报业务中的应用框架构想，强调智能体在异构数据访问、多步骤任务
执行、人机交互三方面的独特优势。陈敏等 \cite{ref9} 在强对流天气
预测的深度学习综述中也表明，气象领域的深度学习应用正向更加智能
化的方向演进。Kim 等 \cite{ref1} 提出的 CLIMATEAGENT 系统进一步
证实了多智能体编排在复杂气象数据科学工作流中的可行性，但其架构
偏向多智能体协作的科研工作流场景，与面向终端用户的轻量化检索
分析需求存在差异。

\subsubsection*{1.2.2 \quad 基于大语言模型的 AI 智能体技术}

AI 智能体是指具备自主感知环境、制定计划、执行行动能力的智能
系统。Xi 等 \cite{ref10} 在综述中系统归纳了基于大语言模型的
智能体的核心能力框架，包含配置 (Profile)、记忆 (Memory)、
规划 (Planning)、行动 (Action) 四个维度。

在推理与规划方向，Wei 等 \cite{ref13} 提出的思维链 (Chain-of-Thought,
CoT) 提示方法显著提升了大语言模型的复杂推理能力，让模型在算术、
常识推理、符号操作任务上的表现接近人类专家水平。Yao 等 \cite{ref30}
提出的思维树 (Tree of Thoughts, ToT) 进一步把推理过程拓展为
树状搜索结构，增强了大语言模型解决需要多步规划问题的能力。最具
影响力的是 Yao 等 \cite{ref18} 提出的 ReAct 范式，该范式把推理
与行动统一在交替循环的框架中：模型先生成思考来分析当前状态，再
决定执行一个动作调用外部工具，最后根据观察到的结果进行下一步
推理。这一范式为构建能够与外部环境交互的智能体提供了基础架构。

在工程框架方向，LangChain \cite{ref15} 提供了构建大语言模型应用
与 Agent 的成熟组件库，支持工具编排、记忆管理与多种 Agent 范式,
是当前应用最广的 Agent 开发框架。Wu 等 \cite{ref19} 提出的
AutoGen 面向多智能体协作场景，让多个具有不同角色的 Agent 通过
结构化对话协作完成复杂任务。Hong 等 \cite{ref20} 的 Data Interpreter
专门面向数据科学任务，通过大语言模型生成并执行代码来完成数据
分析，为解决大语言模型数值计算幻觉问题提供了重要参考。Boiko 等
\cite{ref23} 在自主化学研究中的成功实践展示了大语言模型 Agent
在科学研究领域的潜力。Qian 等 \cite{ref29} 的 ChatDev 工作展示了
通过多智能体协作生成代码的能力，为代码生成驱动的数据分析方案
提供了思路。

\subsubsection*{1.2.3 \quad 检索增强生成与工具学习}

检索增强生成是解决大语言模型知识局限性的核心技术路线。Lewis 等
\cite{ref12} 首次提出 RAG 框架，通过在生成过程中动态检索外部
知识库来增强大语言模型的输出质量。Gao 等 \cite{ref24} 的综述
全面梳理了 RAG 在大语言模型时代的最新进展，涵盖检索策略优化、
知识库构建、端到端训练等方向。黄冰 \cite{ref4} 在古生物学领域
的 RAG 知识问答系统实践表明，RAG 技术能够有效地把领域专业知识
注入到大语言模型中，提升垂直领域的问答准确率。

工具学习是 AI 智能体执行外部操作的关键能力。Schick 等 \cite{ref21}
的 Toolformer 让语言模型学会自主调用计算器、搜索引擎等外部工具,
开创了大语言模型主动使用工具的研究方向。Patil 等 \cite{ref28}
的 Gorilla 专注于大语言模型与大规模 API 的连接，通过检索感知
训练显著提升了 API 调用的准确性。Model Context Protocol
\cite{ref11} 作为新兴的工具协议标准，为智能体与外部工具及数据
源的标准化交互提供了统一规范。林怡等 \cite{ref3} 在 RESTful API
识别方面的工作也为 API 工具封装提供了方法论参考。Gruver 等
\cite{ref22} 的研究表明大语言模型具备零样本时间序列预测能力，
但同时揭示了其在精确数值计算方面的局限性，与 Hong 等 \cite{ref20}
观察到的代码生成必要性形成了相互印证。

\subsubsection*{1.2.4 \quad 现有研究的不足与本课题的定位}

综合上述综述可以归纳出现有研究三方面的不足。

第一是气象 AI 研究偏重预报建模而忽视数据检索与分析的智能化。
现有气象 AI 工作大多致力于提升预报精度，盘古、GraphCast、风乌、
伏羲等大模型 \cite{ref16, ref17, ref26, ref27} 都聚焦于网格化
气象要素的预测精度提升，却缺少面向终端用户的智能数据检索与分析
系统。这导致预报模型产出的高精度数据与终端用户能感知到的决策
价值之间始终存在转化缺口。

第二是通用 Agent 框架缺乏气象领域适配。现有大语言模型 Agent
框架例如 LangChain \cite{ref15} 与 AutoGen \cite{ref19} 虽然
提供了通用的工具调用与多 Agent 协作能力，但未针对气象数据的
异构性、气象术语的专业性、气象分析任务的特殊性进行优化。Kim 等
\cite{ref1} 的 CLIMATEAGENT 是少有的气象领域专门 Agent 工作,
但其多智能体协作架构与面向终端用户的轻量化检索分析需求在工程
形态上存在显著差异。

第三是大语言模型在数值分析中的幻觉问题尚未在气象场景中得到
系统解决。Hong 等 \cite{ref20} 与 Gruver 等 \cite{ref22} 的工作
共同揭示了大语言模型在数值计算与时间序列分析中存在的幻觉，
Data Interpreter \cite{ref20} 提出了通用的代码生成方案，但
气象场景下的数值理解还涉及国家标准引用、分级阈值归属、跨要素
组合判定等更细致的语义层面，需要在通用代码生成方案之上叠加领域
特定的语义桥接机制。如何在气象数据统计分析的具体场景中设计可靠
的确定性计算与权威引用机制，仍缺乏系统性研究。

本课题针对上述三方面不足，提出基于 ReAct 范式的气象 AI 智能体
方案，融合工具学习、RAG 知识注入、代码生成驱动的确定性分析、
语义桥接机制四类技术，构建从自然语言查询到气象决策报告的完整
技术链路。具体研究内容与创新贡献在 \ref{sec:research-content}
节展开。

